% Section 19.1, from exam/README.md and lectures/L20/README.md, with the grading rules of
% info/README.md and the preparation of lectures/L19/README.md.

\section{Inför tentamen}
\label{c19:sec:infor-tentamen}
På tentamen ska du visa att kursens lärandemål är uppnådda. Det här avsnittet beskriver hur den
går till, vad den innehåller och hur du förbereder dig.

\subsection{Format}
\label{c19:sec:format}
Tentamen är en skriftlig salstentamen. Alla uppgifter kräver redovisade lösningar och
motiveringar med angivet svar, inklusive enhet där det är aktuellt. Du måste alltså förklara hur
du kom fram till dina resultat.

\subsection{Tillåtna hjälpmedel}
\label{c19:sec:hjalpmedel}
\begin{itemize}
  \item Skrivmaterial och valfri miniräknare.
  \item En hand- eller datorskriven formelsamling på ett A4-ark (båda sidor).
  \item Ett handskrivet A4-ark (båda sidor) med valfria egna anteckningar.
\end{itemize}

\begin{aside}
\note[Obs] Mobiltelefoner får inte användas under tentamen och ska placeras på angiven plats.
\end{aside}

\subsection{Poäng och betygsgränser}
\label{c19:sec:betyg}
Tentamen ger totalt $25$ poäng. Därtill kan upp till $1{,}0$ bonuspoäng erhållas från kursens två
duggor, $0{,}5$ poäng per dugga för den som får minst $3{,}5$ av duggans $5$ poäng.
Betygsgränserna står i tabellen nedan; för godkänt betyg krävs dessutom att alla kursens
lärandemål är uppnådda.

\begin{narrowtable}{@{}ll@{}}
  \thead{Betyg} & \thead{Gräns} \\
  \midrule
  IG & $< 10{,}0$ poäng \\
  G & $\geq 10{,}0$ poäng \\
  VG & $\geq 17{,}5$ poäng \\
\end{narrowtable}

\subsection{Innehåll}
\label{c19:sec:innehall}
Tentamen är uppdelad i två delar:
\begin{itemize}
  \item \textbf{Del I}: aritmetik, algebra, ekvationer, trigonometri och logaritmer
    (kapitel~\ref{c1:ch}--\ref{c10:ch}).
  \item \textbf{Del II}: derivata, integraler och komplexa tal
    (kapitel~\ref{c12:ch}--\ref{c17:ch}).
\end{itemize}

\subsection{Förberedelse}
\label{c19:sec:forberedelse}
\begin{itemize}
  \item Skriv övningstentamen i \secref{c19:sec:ovningstentamen} på egen hand.
  \item Jämför sedan med facit och med lösningsförslaget, som publiceras i kursens repository, i
    \file{exam/practice_exam_solutions.md}.
  \item Repetera kursens centrala moment, särskilt de områden där du räknade fel på
    övningstentamen.
  \item Förbered formelsamlingen och anteckningsarket (\secref{c19:sec:hjalpmedel}) i god tid.
\end{itemize}
